\documentclass[12pt]{article}
\usepackage[utf8]{inputenc}
\usepackage[T1]{fontenc}
\usepackage[margin=1in]{geometry}
\usepackage{setspace}
\usepackage{xcolor}
\usepackage[colorlinks=true, urlcolor=blue!60!black, linkcolor=blue!60!black]{hyperref}
\setlength{\parindent}{2em}
\setlength{\parskip}{1.25em}
\renewcommand{\baselinestretch}{1.0}
\begin{document}
\begin{center}
{\LARGE \textit{Biden v.\ Nebraska}, 600 U.S.\ 477 (2023)\par}
\vspace{0.8em}
{\large Case Brief\par}
\end{center}
\vspace{1em}

\section{Memory Jogger}

In \href{https://www.law.cornell.edu/supremecourt/text/22-506}{\textit{Biden v. Nebraska}}, the Supreme Court held 6--3 that Missouri had Article III standing via financial harm to its state-created loan servicer MOHELA and that the Secretary of Education's authority under the HEROES Act to ``waive or modify'' student-loan provisions did not authorize a \$430 billion mass student-debt cancellation program under both ordinary statutory text and the Major Questions Doctrine.

\section{Detailed Case Facts}

\textbf{Statutory Framework.} Title IV of the Higher Education Act of 1965 (Education Act), 20 U.S.C. \S~1070(a), governs federal student financial assistance, including Direct Loans (\S~1087a \textit{et seq.}), Perkins Loans (\S~1087aa \textit{et seq.}), and Federal Family Education Loans (FFELs) (\S~1071 \textit{et seq.}). While the Education Act sets loan terms, interest rates, and repayment structures in detail, it authorizes the Secretary of Education to discharge or cancel loan principal only in narrowly enumerated circumstances: qualifying public service (\S\S~1078--10, 1087j, 1087ee), borrower death or permanent and total disability (\S~1087(a)(1)), bankruptcy (\S~1087(b)), and school closure, false certification, or failure to refund loan proceeds (\S~1087(c)).

\textbf{The HEROES Act.} Following the September 11, 2001 terrorist attacks, Congress enacted the Higher Education Relief Opportunities for Students Act of 2001 (115 Stat. 2386) to grant temporary waiver authority for affected borrowers. In 2003, Congress made that authority permanent and extended it to any war, military operation, or national emergency in the \href{https://www.law.cornell.edu/uscode/text/20/chapter-28/subchapter-IV/part-G-1}{HEROES Act of 2003}, Pub. L. No. 108--76, 117 Stat. 904 (codified at 20 U.S.C. \S\S~1098aa--1098ee), passing by a 421--1 vote in the House (with the lone dissenter later stating he meant to vote ``yea'') and a unanimous voice vote in the Senate. Under \href{https://www.law.cornell.edu/uscode/text/20/1098bb}{20 U.S.C. \S~1098bb(a)(1)}, the Secretary ``may waive or modify any statutory or regulatory provision applicable to the student financial assistance programs under title IV'' as the Secretary ``deems necessary in connection with a war or other military operation or national emergency.'' Such waivers or modifications may be issued only ``as may be necessary to ensure'' that ``affected individuals are not placed in a worse position financially in relation to that financial assistance because of their status as affected individuals'' (\S~1098bb(a)(2)(A)). An ``affected individual'' includes anyone who resides or works in a declared disaster area or ``suffered direct economic hardship as a direct result of a war or other military operation or national emergency, as determined by the Secretary'' (\S~1098ee(2)(C)--(D)). Procedurally, \S~1098bb(b)(1)--(2) exempts such actions from APA notice-and-comment rulemaking (5 U.S.C. \S~553) and directs the Secretary to publish a Federal Register notice ``includ[ing] the terms and conditions to be applied in lieu of such statutory and regulatory provisions.''

\textbf{Prior Practice and COVID--19.} Before 2020, HEROES Act waivers and modifications were modest and largely procedural---such as extending filing deadlines, allowing oral authorizations, waiving written leave-of-absence requests, and relaxing uninterrupted-service rules for public-service loan forgiveness (68 Fed. Reg. 69312--69318 (2003); 77 Fed. Reg. 59311--59318 (2012)). On March 13, 2020, President Trump declared the COVID--19 pandemic a national emergency (Proclamation No. 9994). Secretary of Education Betsy DeVos paused loan repayments and interest accrual, Congress extended that pause in the CARES Act through September 2020, and both the Trump and Biden Administrations repeatedly extended the payment pause under the HEROES Act. In January 2021, the Department of Education's Office of General Counsel issued a memorandum concluding that the HEROES Act did not authorize blanket or mass principal cancellation. In August 2022, shortly before President Biden stated that ``the pandemic is over,'' the Office of General Counsel rescinded the 2021 memorandum and concluded that the HEROES Act permitted targeted loan cancellation (87 Fed. Reg. 52944--52945 (2022)).

\textbf{The Debt Cancellation Plan and MOHELA.} In October 2022, Secretary Miguel Cardona published a Federal Register notice (87 Fed. Reg. 61512--61514) establishing a debt cancellation plan for borrowers with eligible federal loans (Direct Loans, Department-held FFELs and Perkins Loans) whose 2020 or 2021 adjusted gross income was below \$125,000 (or \$250,000 for joint filers, heads of household, or qualifying widow(er)s). Qualifying borrowers received up to \$10,000 in principal discharge, or up to \$20,000 if they had received a Pell Grant. With \$1.6 trillion in federal student loans held by 43 million borrowers, the Congressional Budget Office estimated that the plan would cancel roughly \$430 billion in principal (completely erasing debt for 20 million borrowers and cutting the median balance for the other 23 million from \$29,400 to \$13,600), covering an estimated 98.5\% of all borrowers. The Penn Wharton Budget Model projected taxpayer costs between \$469 billion and \$519 billion.

Six States---Nebraska, Missouri, Arkansas, Iowa, Kansas, and South Carolina---sued to block the plan. Missouri's standing theory centered on the Missouri Higher Education Loan Authority (MOHELA), a nonprofit government corporation created by \href{https://revisor.mo.gov/main/PageSelect.aspx?section=173.360&bid=8823}{Mo. Rev. Stat. \S~173.360} as a ``public instrumentality'' performing an ``essential public function.'' MOHELA owned over \$1 billion in FFELs, serviced nearly \$150 billion in federal loans across five million accounts under contract with the Department of Education (generating \$88.9 million in FY 2022 servicing revenue), and used its profits to fund \$230 million in Missouri college development projects and nearly \$300 million in student scholarships. Because discharging half of all federal accounts would cost MOHELA an estimated \$44 million per year in servicing fees, Missouri argued that harm to MOHELA directly injured the State---even though MOHELA declined to join the lawsuit and was legally and financially separate from Missouri's treasury.

\section{Procedural History}

\textbf{U.S. District Court for the Eastern District of Missouri.} On September 29, 2022, the six States filed suit challenging the plan under the separation of powers and the Administrative Procedure Act (APA). On October 20, 2022, Judge Henry Edward Autrey dismissed the action for lack of Article III standing, holding that MOHELA's legal and financial independence prevented Missouri from suing on MOHELA's behalf and rejecting the States' other standing theories (\href{https://clearinghouse.net/case/43874/}{\textit{Nebraska v. Biden}}, 636 F. Supp. 3d 991 (E.D. Mo. 2022)).

\textbf{U.S. Court of Appeals for the Eighth Circuit.} On appeal, the States sought an emergency injunction pending appeal. On November 14, 2022, a unanimous Eighth Circuit panel granted a nationwide preliminary injunction halting the program, holding that Missouri likely had Article III standing through MOHELA, that the merits raised substantial questions, and that the equities favored preserving the status quo (\href{https://law.justia.com/cases/federal/appellate-courts/ca8/22-3179/22-3179-2022-11-14.html}{\textit{Nebraska v. Biden}}, 52 F.4th 1044 (8th Cir. 2022)).

\textbf{U.S. Supreme Court.} On November 18, 2022, the Biden Administration applied to the Supreme Court to vacate the Eighth Circuit's injunction or, alternatively, to grant certiorari before judgment. On December 1, 2022, the Court deferred the vacatur application, construed it as a petition for a writ of certiorari before judgment, granted review on (1) Article III standing and (2) whether the plan exceeded the Secretary's statutory authority or was arbitrary and capricious, and expedited \href{https://www.oyez.org/cases/2022/22-506}{oral argument} for February 28, 2023 alongside a companion challenge by two individual borrowers, \href{https://www.supremecourt.gov/opinions/22pdf/600us1r57_o7kq.pdf}{\textit{Department of Education v. Brown}}, 600 U.S. 551 (2023) (where the Court unanimously held the individual plaintiffs lacked standing). On June 30, 2023, the Supreme Court \textbf{reversed} the District Court's dismissal, \textbf{remanded} for further proceedings, and \textbf{denied as moot} the Government's application to vacate the Eighth Circuit's injunction. On August 16, 2023, the District Court entered final judgment vacating and setting aside the program pursuant to joint stipulation.

\section{Judicial Votes}

\textbf{Vote:} 6--3 reversing the District Court's dismissal and invalidating the debt cancellation program (\href{https://www.supremecourt.gov/opinions/22pdf/600us1r56_1o13.pdf}{600 U.S. 477 (2023)}).

\begin{itemize}
  \item \textbf{Majority Opinion:} Chief Justice Roberts delivered the opinion of the Court, joined in full by Justices Thomas, Alito, Gorsuch, Kavanaugh, and Barrett.
  \item \textbf{Concurring Opinion:} Justice Barrett filed a sole concurring opinion (joining the majority in full) to defend the Major Questions Doctrine as a contextual linguistic tool consistent with textualism rather than a substantive clear-statement canon.
  \item \textbf{Dissenting Opinion:} Justice Kagan filed a dissenting opinion, joined by Justices Sotomayor and Jackson.
\end{itemize}

\section{Holding}

\textbf{Judgment.} Reversed and remanded (reversing the Eastern District of Missouri's dismissal and invalidating the Secretary of Education's student-loan cancellation plan); the Government's application to vacate the Eighth Circuit's injunction was denied as moot.

\textbf{Rules Announced.}
\begin{enumerate}
  \item \textbf{Article III Standing for State Instrumentalities:} When a State creates, governs, supervises, and retains power to dissolve a public corporation performing an essential public function (an ``instrumentality'' of the State), a direct, concrete financial injury to that instrumentality in performing its public mission constitutes a direct injury in fact to the State itself, conferring Article III standing even if the corporation has a separate legal personality and can sue in its own name.
  \item \textbf{Scope of ``Waive or Modify'' under the HEROES Act:} Authority under \href{https://www.law.cornell.edu/uscode/text/20/1098bb}{20 U.S.C. \S~1098bb(a)(1)} to ``waive or modify'' existing statutory or regulatory provisions permits only nullification of specific requirements or modest, incremental adjustments---not the wholesale creation of a novel, \$430 billion loan-forgiveness regime replacing Congress's carefully delineated statutory scheme.
  \item \textbf{Major Questions Doctrine Applied to Benefits Programs:} Separation-of-powers concerns under the Major Questions Doctrine apply equally to federal spending and benefits programs as to regulatory mandates; where an agency asserts unheralded authority of vast economic and political significance, it must point to ``clear congressional authorization.''
\end{enumerate}

\textbf{Holding vs. Dicta.} The holdings are (1) Missouri's standing via MOHELA and (2) the invalidity of the debt cancellation plan under the HEROES Act's text and the Major Questions Doctrine. Several passages are \textit{dicta}: (a) Footnote 5 explicitly declines to decide the legality of the COVID--19 payment and interest-accrual pauses implemented under the HEROES Act; (b) Footnote 6 notes in passing---without resting the decision on it---that the Secretary would face a ``daunting task'' showing mass cancellation was ``necessary'' under \S~1098bb(a)(2)(A) given the prior payment pause; (c) the Court declines to address the other five States' standing theories; and (d) Justice Barrett's lone concurrence reframing the Major Questions Doctrine as a linguistic canon rather than a substantive clear-statement rule represents her individual methodology rather than binding Court precedent.

\section{Analysis of Opinions}

\subsection{Majority Opinion (Chief Justice Roberts)}

\textbf{Article III Standing (Part II).} Applying \href{https://supreme.justia.com/cases/federal/us/504/555/}{\textit{Lujan v. Defenders of Wildlife}}, 504 U.S. 555, 560--561 (1992), Roberts noted that only one plaintiff needs standing for a suit to proceed. Both the Government and the dissent conceded that discharging half of federal loan accounts would cost MOHELA roughly \$44 million annually in servicing fees---a classic financial injury directly traceable to the plan. The dispositive question was whether MOHELA's injury was Missouri's injury. Roberts cataloged Missouri's control over MOHELA: Missouri created it by statute (\href{https://revisor.mo.gov/main/PageSelect.aspx?section=173.360&bid=8823}{Mo. Rev. Stat. \S~173.360}) as a ``public instrumentality'' to perform the ``essential public function'' of expanding higher-education access; its board consists of two state officials and five gubernatorial appointees confirmed by the state Senate and removable for cause; it submits annual financial reports to the State; and only Missouri can dissolve it. Relying on \href{https://supreme.justia.com/cases/federal/us/346/368/}{\textit{Arkansas v. Texas}}, 346 U.S. 368 (1953) (allowing Arkansas to sue over contract interference harming the University of Arkansas), \href{https://supreme.justia.com/cases/federal/us/513/374/}{\textit{Lebron v. National Railroad Passenger Corp.}}, 513 U.S. 374, 397 (1995) (treating Amtrak as part of the Government for First Amendment purposes despite corporate separateness), and \textit{Department of Transportation v. Association of American Railroads}, 575 U.S. 43 (2015), Roberts concluded that a State's choice to carry out public functions through a separate corporate entity does not bar the State from suing to remedy direct harm to that instrumentality.

\textbf{Textual Analysis of ``Waive or Modify'' (Part III--A).} On the merits, Roberts began with ordinary statutory interpretation. Drawing heavily on Justice Scalia's opinion in \href{https://supreme.justia.com/cases/federal/us/512/218/}{\textit{MCI Telecommunications Corp. v. AT\&T Co.}}, 512 U.S. 218, 225, 228 (1994), Roberts held that ``modify'' carries a connotation of increment or limitation---``to change moderately or in minor fashion''---and cannot authorize ``basic and fundamental changes in the scheme'' designed by Congress. Transforming narrow discharge exceptions (death, disability, bankruptcy, school closure) into automatic cancellation for 98.5\% of borrowers ``modified'' the Education Act ``only in the same sense that `the French Revolution ``modified'' the status of the French nobility.''' Nor could ``waive'' save the plan: the Secretary did not merely nullify a discrete statutory requirement (as in prior HEROES Act waivers) because no Education Act provision imposes the repayment obligation itself; instead, the Secretary deleted existing statutory limits and substituted an elaborate new regime of dollar caps and income thresholds. Rejecting the Government's argument that ``waive or modify'' is greater than the sum of its parts, Roberts also dismissed \S~1098bb(b)(2)'s reference to publishing ``terms and conditions to be applied in lieu of'' waived provisions as a ``humdrum reporting requirement'' and ``a wafer-thin reed'' that remains bounded by the incremental scope of ``modify.''

\textbf{Major Questions Doctrine (Part III--B).} Reinforcing the textual holding, Roberts invoked \href{https://supreme.justia.com/cases/federal/us/597/20-1530/}{\textit{West Virginia v. EPA}}, 597 U.S. 697 (2022), \href{https://supreme.justia.com/cases/federal/us/529/120/}{\textit{FDA v. Brown \& Williamson Tobacco Corp.}}, 529 U.S. 120, 159--160 (2000), and \href{https://supreme.justia.com/cases/federal/us/573/302/}{\textit{Utility Air Regulatory Group v. EPA}}, 573 U.S. 302, 324 (2014). Four indicators triggered heightened skepticism: (1) \textit{Unheralded power}: no prior HEROES Act action approached this scale; (2) \textit{Staggering economic significance}: \$430 billion to \$519 billion equals nearly one-third of annual federal discretionary spending and ten times the economic impact of the CDC eviction moratorium struck down in \textit{Alabama Assn. of Realtors v. HHS}, 594 U.S. 758 (2021); (3) \textit{Political salience and congressional inaction}: Congress had considered over 80 student-loan forgiveness bills without enacting mass cancellation, and then-Speaker Nancy Pelosi had publicly stated in July 2021 that the President lacked authority to forgive debt; and (4) \textit{Control of the purse}: rejecting the Government's distinction between regulatory burdens and government benefits, Roberts cited the Appropriations Clause (Art. I, \S~9, cl. 7) and \href{https://supreme.justia.com/cases/federal/us/576/473/}{\textit{King v. Burwell}}, 576 U.S. 473, 485 (2015), to confirm that Major Questions principles apply to massive spending and benefits programs. Requiring ``clear congressional authorization,'' the Court found none.

\subsection{Concurring Opinion (Justice Barrett)}

Joining the majority in full, Justice Barrett wrote separately to reconcile the Major Questions Doctrine with textualism. She candidly acknowledged that ``strong-form substantive canons''---which impose a ``clarity tax'' forcing courts to reject the best reading of a statute in favor of an inferior-but-tenable reading to protect external constitutional values---are ``in significant tension with textualism.'' Rather than viewing the Major Questions Doctrine as a substantive nondelegation clear-statement rule (a framing associated with Justice Gorsuch's concurrence in \textit{West Virginia}), Barrett reframed it as a \textit{linguistic and contextual} tool of ordinary statutory interpretation.

Drawing on agency law (Restatement (Third) of Agency \S~2.02) and everyday communication, Barrett offered her now-famous \textbf{babysitter analogy}: if a parent hands a babysitter a credit card and says, ``Make sure the kids have fun,'' literalism might allow a two-day, out-of-town amusement-park road trip with a hotel stay, yet common sense and context (including who holds primary authority) show that a reasonable speaker would expect explicit authorization for such an extraordinary excursion. Likewise, because Article I, \S~1 vests ``[a]ll legislative Powers'' in Congress, constitutional structure is part of the background context against which reasonable interpreters read delegations. Context---such as an agency acting outside its ``wheelhouse'' (as in \href{https://supreme.justia.com/cases/federal/us/595/21a244/}{\textit{NFIB v. OSHA}}, 595 U.S. 109 (2022), or \textit{Gonzales v. Oregon}, 546 U.S. 243 (2006)), relying on a ``backwater'' provision (\textit{West Virginia}; \href{https://supreme.justia.com/cases/federal/us/531/457/}{\textit{Whitman v. American Trucking Assns.}}, 531 U.S. 457, 468 (2001)), or claiming unheralded economic power---is not an ``on-off switch'' but a continuum of clues. Even though the Secretary of Education was not operating wholly outside his usual domain here, the remaining contextual indicators confirmed that the best reading of the HEROES Act excluded \$430 billion in mass debt cancellation.

\subsection{Dissenting Opinion (Justice Kagan)}

Joined by Justices Sotomayor and Jackson, Justice Kagan delivered a forceful dissent attacking both standing and the merits as judicial overreach.

\textbf{Standing (Part I).} Kagan argued that the six States were ``classic ideological plaintiffs'' with no concrete stake in the program. Missouri's treasury would not lose ``one penny'' because MOHELA is legally and financially independent under Missouri law: MOHELA's assets are not state revenue (Mo. Rev. Stat. \S\S~173.386, 173.425), Missouri is not liable for MOHELA's debts (\S~173.410), and MOHELA has independent power to ``sue and be sued'' (\S~173.385.1(3)). Asking ``Where's MOHELA?,'' Kagan emphasized that MOHELA declined to join the suit, filed no amicus brief, and forced the Missouri Attorney General to obtain its records through formal state sunshine-law requests. She cited \textit{Menorah Medical Center v. Health and Ed. Facilities Auth.}, 584 S.W.2d 73, 78 (Mo. 1979), where the Missouri Supreme Court held an analogous state authority was ``not the State.'' Kagan distinguished \href{https://supreme.justia.com/cases/federal/us/346/368/}{\textit{Arkansas v. Texas}} (where the State owned all university property and the Court believed the university could not sue on its own) and \href{https://supreme.justia.com/cases/federal/us/513/374/}{\textit{Lebron}} (which expressly limited its state-actor holding to constitutional rights and reaffirmed that a sovereign is ``not a party to suits brought by or against'' its corporation). Moreover, invoking harm to Missouri's college students violated the bar on state \textit{parens patriae} suits against the Federal Government (\textit{Massachusetts v. Mellon}, 262 U.S. 447 (1923)). Turning Chief Justice Roberts's own dissent in \textit{Massachusetts v. EPA}, 549 U.S. 497, 548 (2007), against him, Kagan accused the majority of turning standing into ``a lawyer's game.''

\textbf{Merits and Major Questions (Part II).} On the text, Kagan accused the majority of ``divide to conquer'' textualism---sealing ``waive,'' ``modify,'' and ``in lieu of'' in separate vacuum-packed containers rather than reading the whole sentence together. Because ``waive'' permits wholly eliminating a requirement (the largest possible change) and ``modify'' permits altering it, the couplet ``waive or modify'' naturally authorizes changes along a spectrum from minor to substantial, and \S~1098bb(b)(2)'s ``in lieu of'' clause confirms the Secretary may substitute new terms so waived rules do not leave ``gaping holes.'' Kagan offered a \textbf{Chicago dirty-bomb hypothetical}: if terrorists detonated a dirty bomb in Chicago, forcing borrowers to flee, the HEROES Act was unmistakably designed to let the Secretary cancel \$10,000 of their loans without waiting for Congress; under the majority's cramped reading, the Secretary could ``do no more than fiddle'' with tax forms, and even the \$100+ billion COVID--19 payment pause would be unlawful. Finally, Kagan blasted the majority's invocation of the ``made-up major-questions doctrine'' as a ``tell'' that ordinary text could not sustain the outcome. Noting in Footnote 3 that under Justice Barrett's own contextual factors there was neither a ``backwater'' provision nor an agency acting outside its ``wheelhouse,'' Kagan argued that the Court had moved the goalposts to veto broad emergency delegations whenever it disliked their scale, substituting unelected judges for democratically accountable branches.

\section{Examples: Future Applications}

\begin{enumerate}
  \item \textbf{Same Side (Invalid Agency Action): Unilateral Agricultural Debt Cancellation Under a Disaster ``Modify'' Clause.} Suppose the Secretary of Agriculture invokes a post-9/11 farm-relief statute authorizing the Department to ``waive or modify any statutory or regulatory provision'' governing federal farm loans ``as necessary in connection with a national emergency'' to cancel \$350 billion in principal across 95\% of U.S. farms during an inflationary spike declared an economic emergency. Under \textit{Biden v. Nebraska}, a state-chartered agricultural lending authority losing servicing fees would confer standing on its parent State, and the program would fail both textually (creating a novel forgiveness regime exceeds ``waive or modify'') and under the Major Questions Doctrine (unheralded power of vast fiscal magnitude implicating Congress's Appropriations Clause authority).

  \item \textbf{Same Side (Invalid Agency Action): Nationwide Mortgage Principal Reduction by FHFA/HUD.} Suppose HUD invokes an emergency housing statute permitting it to ``waive or modify'' FHA mortgage insurance rules during a presidentially declared housing emergency to forgive \$50,000 of principal on all federally backed mortgages for households earning under \$150,000, at a projected \$500 billion cost to the Treasury. Even though housing finance sits within HUD's general domain, \textit{Biden v. Nebraska} forecloses using ``waive or modify'' to invent a sweeping debt-cancellation program where Congress has enacted only narrow hardship-forbearance provisions and repeatedly declined to pass mass mortgage-forgiveness bills.

  \item \textbf{Opposite Side (Valid Agency Action): Targeted Procedural Waivers After a Hurricane.} Suppose after a Category 5 hurricane triggers a presidential national emergency declaration in coastal Texas and Louisiana, the Secretary of Education invokes the HEROES Act to waive written-verification deadlines, extend loan-rehabilitation windows by 12 months, and relax uninterrupted-service requirements for Public Service Loan Forgiveness strictly for borrowers residing in FEMA-designated disaster counties. This closely mirrors the modest, targeted 2003 and 2012 waivers that Chief Justice Roberts explicitly endorsed in Part III--A as faithful exercises of ``waive or modify,'' and it triggers none of the economic, political, or anti-novelty red flags of the Major Questions Doctrine.

  \item \textbf{Opposite Side (Valid Agency Action): Clear Statutory Authorization for Means-Tested Relief.} Suppose Congress amends the Higher Education Act to enact a new provision expressly stating: ``During any national emergency declared under the National Emergencies Act, the Secretary of Education may cancel or discharge up to \$15,000 in principal per borrower for borrowers with adjusted gross income below \$100,000, up to an aggregate fiscal cap of \$200 billion.'' Because Congress has supplied ``clear congressional authorization'' specifying the mechanism, caps, and fiscal ceiling, a subsequent cancellation plan within those bounds would satisfy both ordinary textualism and the Major Questions Doctrine under both Roberts's majority opinion and Barrett's contextual framework.

  \item \textbf{Fence-Sitter (Close Case): Indefinite National Payment and Interest Pause Under the HEROES Act.} Suppose during a severe, multi-year pandemic, the Secretary of Education does not cancel any loan principal, but instead invokes the HEROES Act to suspend all federal student-loan payments and set interest rates to 0\% for three years, costing the federal government \$150 billion in forgone interest while leaving statutory discharge provisions untouched. Chief Justice Roberts expressly left the legality of the COVID--19 payment pause undecided in Footnote 5. On the Government's side, setting payments and interest to zero looks closer to a temporary ``waiver'' of repayment timing and accrual than inventing a permanent principal-cancellation regime, and Congress implicitly ratified the pause in the CARES Act. On the challengers' side, as Justice Kagan pointed out in dissent, a \$150 billion nationwide forbearance for 43 million borrowers is hardly a ``minor'' or ``incremental'' adjustment under \textit{MCI} and rivals the economic magnitude that triggered the Major Questions Doctrine in \textit{Nebraska}.
\end{enumerate}

\section{Critique}

\textbf{Progressive and Institutionalist Critiques.} Progressive scholars and administrative-law commentators attacked \textit{Biden v. Nebraska} on both Article III and interpretive grounds. On standing, critics echoed Justice Kagan's warning that letting Missouri sue on behalf of an unwilling, financially insulated public corporation eroded Article III's injury-in-fact guardrail and enabled ideological forum shopping by state attorneys general. Strikingly, subsequent federal appellate litigation exposed the doctrinal tension in Roberts's ``arm of the State'' reasoning: in \href{https://law.justia.com/cases/federal/appellate-courts/ca10/22-3286/22-3286-2024-11-12.html}{\textit{Good v. U.S. Department of Education}}, 121 F.4th 772 (10th Cir. 2024), the Tenth Circuit held that MOHELA is \textit{not} an ``arm of the State'' entitled to Eleventh Amendment sovereign immunity precisely because Missouri's treasury is insulated from MOHELA's liabilities---producing the asymmetry that MOHELA counts as the State when Missouri wants to sue, but not when MOHELA is sued. On the merits, scholars such as \href{https://www.cato.org/sites/cato.org/files/2023-09/cato-supreme-court-review-9.pdf}{Jed Handelsman Shugerman} (2022--2023 Cato Sup. Ct. Rev. 209) argued that Roberts's opinion relied on an old-fashioned purposivist Major Questions Doctrine that effectively rewrote broad statutory text (``waive or modify \textit{any} statutory or regulatory provision'') to match judicial intuitions about what Congress ``would likely have intended.''

\textbf{Empirical and Linguistic Critiques of Barrett's Concurrence.} Justice Barrett's effort to rescue the Major Questions Doctrine as a ``common sense'' linguistic canon drew sharp empirical and theoretical scrutiny. In \href{https://www.law.nyu.edu/sites/default/files/tobia_walters_slocum_major\%20questions.pdf}{\textit{Major Questions, Common Sense?}}, Kevin Tobia, Daniel E. Walters, and Brian Slocum empirically tested Barrett's babysitter and grocer-clerk analogies on ordinary Americans and found that respondents did not uniformly treat high-stakes actions as excluded by open-ended delegations---undermining the claim that the doctrine merely mirrors ordinary conversational pragmatics. Similarly, \href{https://www.yalejreg.com/nc/lets-talk-about-that-barrett-concurrence-on-the-contextual-major-questions-doctrine-by-beau-j-baumann/}{Beau J. Baumann} noted that Barrett's contextual approach still smuggles normative separation-of-powers assumptions into ``ordinary meaning'' without transparent limiting principles.

\textbf{Originalist, Textualist, and Conservative Critiques.} Criticism was not confined to the left; prominent conservative, originalist, and textualist scholars also faulted key aspects of the decision. On standing, \href{https://harvardlawreview.org/print/vol-137/proper-parties-proper-relief/}{William Baude and Samuel L. Bray} (137 Harv. L. Rev. 153 (2023)) highlighted the ``state-as-a-super-big-person'' problem, observing that although the Court quietly buried \textit{Massachusetts v. EPA}'s ``special solicitude'' for States, accepting Missouri's borderline MOHELA theory still stretched traditional proper-party limits. On statutory interpretation, conservative textualists such as \href{https://journals.law.harvard.edu/jlpp/placing-legal-context-in-context-chad-squitieri/}{Chad Squitieri} (Harv. J.L. \& Pub. Pol'y Per Curiam (2024)) and \href{https://www.yalejreg.com/nc/text-and-context-by-adrian-vermeule/}{Adrian Vermeule} (Yale J. on Reg. Notice \& Comment (2023)) argued that Barrett's ``contextual'' Major Questions Doctrine risks collapsing textualism into amorphous purposivism or classic-law reasonableness, while strict nondelegation originalists contend that if the HEROES Act truly delegated unbounded power to cancel \$430 billion in debt, the Court should enforce Article I's Vesting Clause directly via a revitalized nondelegation doctrine rather than straining statutory text or inventing ad hoc clear-statement rules.

\textbf{Competing Constitutional Values.} Ultimately, \textit{Biden v. Nebraska} crystallizes an enduring clash between two constitutional visions: a formalist separation-of-powers model prioritizing Congress's exclusive power of the purse and legislative accountability over major fiscal redistributions, versus a functionalist administrative-state model emphasizing Congress's pragmatic need to delegate flexible emergency powers to expert executive agencies accountable through the President.

\section{Key Quotations}

\begin{quote}
``By law and function, MOHELA is an instrumentality of Missouri: It was created by the State to further a public purpose, is governed by state officials and state appointees, reports to the State, and may be dissolved by the State. The Secretary's plan will cut MOHELA's revenues, impairing its efforts to aid Missouri college students. This acknowledged harm to MOHELA in the performance of its public function is necessarily a direct injury to Missouri itself.''
\end{quote}
\noindent --- Chief Justice Roberts, Opinion of the Court, \href{https://www.law.cornell.edu/supremecourt/text/22-506}{600 U.S. 477 (2023)}.

\begin{quote}
``We hold today that the Act allows the Secretary to `waive or modify' existing statutory or regulatory provisions applicable to financial assistance programs under the Education Act, not to rewrite that statute from the ground up.''
\end{quote}
\noindent --- Chief Justice Roberts, Opinion of the Court, \href{https://www.law.cornell.edu/supremecourt/text/22-506}{600 U.S. 477 (2023)}.

\begin{quote}
``The dissent is correct that this is a case about one branch of government arrogating to itself power belonging to another. But it is the Executive seizing the power of the Legislature.''
\end{quote}
\noindent --- Chief Justice Roberts, Opinion of the Court, \href{https://www.law.cornell.edu/supremecourt/text/22-506}{600 U.S. 477 (2023)}.

\begin{quote}
``Rather, I understand it to emphasize the importance of context when a court interprets a delegation to an administrative agency. Seen in this light, the major questions doctrine is a tool for discerning---not departing from---the text's most natural interpretation.''
\end{quote}
\noindent --- Justice Barrett, Concurring Opinion, \href{https://www.law.cornell.edu/supremecourt/text/22-506}{600 U.S. 477 (2023)}.

\begin{quote}
``So in a case not a case, the majority overrides the combined judgment of the Legislative and Executive Branches, with the consequence of eliminating loan forgiveness for 43 million Americans. I respectfully dissent from that decision.''
\end{quote}
\noindent --- Justice Kagan, Dissenting Opinion, \href{https://www.law.cornell.edu/supremecourt/text/22-506}{600 U.S. 477 (2023)}.


\end{document}
